\chapter*{Введение}
Задача автоматической оценки возраста человека по изображению лица является актуальной в современных системах компьютерного зрения и находит применение в задачах биометрической идентификации, анализа пользовательской аудитории и интеллектуальных интерфейсов. Сложность данной задачи обусловлена высокой вариативностью внешности людей, зависимостью внешних признаков от условий съёмки, а также наличием шума и дисбаланса в данных.

Целью данной курсовой работы является разработка и исследование модели оценки возраста человека по изображению лица на основе методов глубокого обучения, в частности сверточных нейронных сетей.

Для достижения поставленной цели в работе были сформулированы и решены следующие подзадачи:
\begin{enumerate}
	\item выполнить сбор и предобработку изображений лиц людей;
	\item реализовать модель машинного обучения;
	\item обучить модель с использованием обучающей и валидационной выборок;
	\item выполнить оценку качества модели и интерпретировать полученные результаты;
\end{enumerate}


\chapter{Постановка задачи}
В рамках работы необходимо реализовать приложение для распознавания возраста по фотографии используя CNN (convolutional neural network).


\textbf{Входные данные:}
\begin{itemize}[]
	\item $F$ - изображение лица человека.
\end{itemize}


\textbf{Выходные данные:}
\begin{itemize}[]
	\item $a \in [0;116]$ - возраст человека предсказанный по фотографии лица.
\end{itemize}

\textbf{Ограничения:}
\begin{itemize}[]
	\item изображения должны содержать одно лицо;
	\item размер изображения 224 на 224 пикселя;
	\item предсказание выполняется только для людей;
	\item в расчет не берутся этнические и половые признаки.
\end{itemize}